\documentclass[12pt]{article}

% --- Paquetes Esenciales ---
\usepackage[utf8]{inputenc}    % Soporte para caracteres en español
\usepackage[spanish]{babel}   % Idioma español (títulos como Figura, Resumen, etc.)
\usepackage{graphicx}         % Para insertar imágenes
\usepackage{amsmath}          % Paquete matemático avanzado
\usepackage{geometry}         % Configuración de márgenes estándar
\geometry{letterpaper, margin=1in}

% --- Datos del Documento ---
\title{Práctica en \LaTeX}
\author{Brandon Alexis Avalos Santillán}
\date{11 de febrero de 2026}

\begin{document}

\maketitle

\section{Introducción}
\LaTeX{} es una herramienta utilizada para la elaboración de documentos técnicos y científicos. En la ingeniería mecatrónica se usa frecuentemente para reportes, artículos y documentación de proyectos.

\section{¿Por qué usar \LaTeX?}
Algunas ventajas de \LaTeX{} son:
\begin{itemize}
    \item Excelente manejo de ecuaciones.
    \item Documentos con formato profesional.
    \item Ideal para trabajos de ingeniería.
\end{itemize}

\section{Ejemplo de ecuación}
Una ecuación común en mecatrónica es la segunda ley de Newton:
\begin{equation}
F = m \cdot a
\end{equation}

Otra ecuación utilizada es la ley de Ohm, la cual relaciona la corriente con el voltaje y la resistencia en un circuito eléctrico:
\begin{equation}
I = \frac{V}{R}
\end{equation}

Donde $I$ es la corriente eléctrica, $V$ es el voltaje y $R$ es la resistencia. Se miden en amperios (A), voltios (V) y ohmios ($\Omega$).

\section{Uso de imágenes}
Para importar imágenes en \LaTeX, asegúrate de subir el archivo al proyecto de Overleaf.

\begin{figure}[h]
    \centering
    % Nota: Asegúrate de tener una imagen llamada "i1.jpg" o con otra extensión válida en tu proyecto de Overleaf
    \includegraphics[width=0.5\textwidth]{i1.jpg} 
    \caption{Imagen 1}
    \label{fig:imagen1}
\end{figure}

\section{Conclusión}
\LaTeX{} es una herramienta poderosa que nos ayuda a presentar nuestros trabajos de forma clara y profesional.

\end{document}